\documentclass[border=10pt]{standalone}
\usepackage{fontspec}
\AtBeginDocument{\hyphenation{zenodo}}% never hyphenate "zenodo" inside a DOI
\setmainfont{Noto Serif}[Ligatures=TeXOff]
\XeTeXgenerateactualtext=1
\newfontfamily\bengalifont{Noto Serif Bengali}[Script=Bengali, NFSSFamily=bnfam]
\usepackage[Latin,Bengali]{ucharclasses}
% group-free font switching (as in header-common.tex): robust inside braces, cells and quotes
\makeatletter
\def\BN@saved{\rmdefault}
\def\BN@bn{bnfam}
\newcommand\BNin{\edef\BN@tmp{\f@family}\ifx\BN@tmp\BN@bn\else\let\BN@saved\BN@tmp\fi\fontfamily{bnfam}\selectfont}
\newcommand\BNout{\fontfamily{\BN@saved}\selectfont}
\makeatother
\setTransitionsFor{Bengali}{\BNin}{\BNout}
% Typographic quotes and dashes are not in a ucharclasses class (class 0), so a Bangla word right
% after “ or ‘ would not trigger the Bangla font. Treat them as Latin so that every switch fires.
\XeTeXcharclass"2018=\XeTeXcharclass`A
\XeTeXcharclass"2019=\XeTeXcharclass`A
\XeTeXcharclass"201C=\XeTeXcharclass`A
\XeTeXcharclass"201D=\XeTeXcharclass`A
\XeTeXcharclass"2013=\XeTeXcharclass`A
\XeTeXcharclass"2014=\XeTeXcharclass`A
\XeTeXcharclass"2026=\XeTeXcharclass`A
\usepackage{newunicodechar}
\newunicodechar{।}{{\bengalifont\char"0964}}
\newunicodechar{॥}{{\bengalifont\char"0965}}

\usepackage[table]{xcolor}
\usepackage{booktabs,tabularx,array}
\definecolor{bapilrule}{HTML}{1F2937}
\definecolor{bapilstripe}{HTML}{F3F4F6}
\renewcommand{\arraystretch}{1.28}
\setlength{\tabcolsep}{5pt}

\begin{document}
\begin{minipage}{17cm}
{\large\bfseries সারণি ৬২।}\ {\large ব্যয়ের ধরনভিত্তিক বাজেট, নিম্ন, মূল ও উচ্চ পরিস্থিতিসহ (কোটি টাকা)}\par\vspace{7pt}
\small
\rowcolors{2}{bapilstripe}{white}
\begin{tabularx}{\linewidth}{@{}>{\hsize=2.530\hsize\raggedright\arraybackslash}X>{\hsize=0.602\hsize\raggedright\arraybackslash}X>{\hsize=0.602\hsize\raggedright\arraybackslash}X>{\hsize=0.602\hsize\raggedright\arraybackslash}X>{\hsize=0.663\hsize\raggedright\arraybackslash}X@{}}
\arrayrulecolor{bapilrule}\toprule
\rowcolor{white}\textbf{বিষয়} & \textbf{বছর ১} & \textbf{বছর ২} & \textbf{বছর ৩} & \textbf{মোট} \\
\midrule
এককালীন প্রত্যক্ষ ব্যয় (যন্ত্রপাতি, বেঞ্চমার্ক তৈরি, সমীক্ষা ও মূল্যায়ন) & ১.৪২ & ১.৪৬ & ০.৭৯ & ৩.৬৭ \\
এর মধ্যে মূলধনি ব্যয় (সার্ভার, GPU ও স্টোরেজ) & ১.১৪ & ১.০২ & ০.৩০ & ২.৪৬ \\
পুনরাবৃত্ত প্রত্যক্ষ ব্যয় (জনবল, সম্পাদক, ফেলো, হোস্টিং, OCR, প্রশিক্ষণ, আশ্বাস) & ২.৬৯ & ৫.০৫ & ৬.৩৫ & ১৪.১০ \\
স্বাগতিকের পরোক্ষ ব্যয় ও আপৎকালীন সংস্থান & ০.৬২ & ০.৯৮ & ১.০৭ & ২.৬৬ \\
\textbf{মূল হিসাবের মোট} & \textbf{৪.৭৩} & \textbf{৭.৪৯} & \textbf{৮.২১} & \textbf{২০.৪৩} \\
নিম্ন পরিস্থিতির মোট & ৪.২০ & ৬.৫৮ & ৭.১৬ & ১৭.৯৩ \\
উচ্চ পরিস্থিতির মোট & ৫.৭৬ & ৯.৪৬ & ১০.০৩ & ২৫.২৫ \\
সমন্বিত মূল হিসাব: জানা বাদ-পড়া ব্যয় যুক্ত ও ১০\% আপৎকালীন সংস্থান (অনুচ্ছেদ ১৭.২) & ৫.২৯ & ৮.১০ & ৮.৮৫ & ২২.২৪ \\
\bottomrule
\end{tabularx}
\par\vspace{5pt}{\footnotesize\itshape নিম্ন পরিস্থিতি: জনবল-ব্যয় ১৫\% কম, মূল্যবৃদ্ধি ৬\%, যন্ত্রপাতি ২০\% সস্তা, প্রতি মার্কিন ডলারে ১২০ টাকা। উচ্চ পরিস্থিতি: জনবল-ব্যয় ২০\% বেশি, মূল্যবৃদ্ধি ১০\%, যন্ত্রপাতি ২৫\% বেশি দামি (যেমন আমদানি শুল্ক ও ভ্যাট), প্রতি মার্কিন ডলারে ১৩৫ টাকা, আপৎকালীন সংস্থান ১০\% এবং পর্যায় ১ অভিযোজনের জন্য ২০,০০০ মার্কিন ডলারের ক্লাউড কম্পিউটিং। সমন্বিত মূল হিসাব: মূল হিসাবের অনুমান, সঙ্গে সম্পৃক্ততা-অনুষ্ঠান, P5-এর জরিপ-তথ্য, ভ্রমণ ও উন্মুক্ত-প্রবেশাধিকার ফির জন্য বছরে ১৫ লাখ টাকা, ১৯ জন নতুন কর্মীর প্রত্যেকের জন্য একটি কম্পিউটার (১,২০০ মার্কিন ডলার) এবং ১০\% আপৎকালীন সংস্থান। প্রতিটি অঙ্ক আলাদাভাবে আসন্নীকৃত।}
\par\vspace{3pt}{\scriptsize\color{gray} উৎস: BAPIL টেকনিক্যাল ব্লুপ্রিন্ট, সংস্করণ ১.০ (২০২৬)। CC BY 4.0।}
\end{minipage}
\end{document}
